\section{门电路}
\subsection{半导体二极管门电路}
\subsubsection{半导体二极管的开关特性}
\textbf{稳态开关特性}\quad 将二极管视作正向导通开启电压 $V_{ON}=0.7 V$, 正向导通阻值为 0, 反向电阻无穷大.

\textbf{动态特性}\quad 二极管两端电压突然反向时:
\begin{itemize}
    \item \underline{外加电压由反向变为正向}: 正向导通电流的建立稍微滞后;
    \item \underline{外加电压由正向变为反向}: 有较大瞬态反向电流流过, 后迅速衰减并趋向饱和电流, 需要一段时间.
\end{itemize}

在输入频率较低时可以忽略, 高频时不可忽略.

\subsubsection{二极管与门}
\subsubsection{二极管或门}
\subsubsection{二极管门电路的缺点}
\begin{itemize}
    \item 电平有偏移 (输出与输入相差二极管压降; 后级门电路电压偏差值会超过门限);
    \item 带负载能力差.
\end{itemize}
